\documentclass[UTF8,11pt]{ctexart}
\usepackage[a4paper,margin=2.5cm]{geometry}
\usepackage{amsmath,amssymb,bm}
\usepackage{booktabs,longtable,array}
\usepackage{enumitem}
\usepackage{hyperref}
\hypersetup{hidelinks}
\setlength{\parindent}{2em}
\setlength{\parskip}{0.35em}
\renewcommand{\arraystretch}{1.25}
\setlist{itemsep=0.2em,topsep=0.3em}
\newcommand{\pos}[1]{\left[#1\right]^+}
\newcommand{\E}{\mathbb E}
\newcommand{\R}{\mathbb R}
\newcommand{\CVaR}{\operatorname{CVaR}}
\allowdisplaybreaks[2]
\title{问题3的详细数学模型\\
{\large 预测提前量可信度、联合场景与CVaR随机滚动调度}}
\author{}
\date{}
\begin{document}
\maketitle

\noindent\textbf{文档性质。}
本文是供团队推导与后续实现的建模说明，不是参赛论文正文，不含求解代码、
附件计算结果或未经验证的节省率。延续问题1、2的中文排版及电量变量口径。
按当前要求，不建立鲁棒优化、分布鲁棒优化或Tube MPC模型。
主方案为\textbf{按提前量校准的概率预测与随机MPC}，
CVaR是与风险中性版本比较的风险增强项。
文献方法与本题提出的指数形式分别标明，不将组合方法直接宣称为学术首创。

\section{任务、物理关系与建模主线}
\subsection{第三问比第二问增加了什么}
每天0:00制定当天的计划购电量；6:00、12:00、18:00取得新光伏预报后，
在支付调整费用的条件下修改尚未执行的购电量，重新安排电池充放电；
实际仍不足的部分以5倍电价紧急购买。
目标是在满足负荷与储能安全约束的前提下，权衡正常购电、调整和紧急费用，
并分析新增预报时刻的价值。

\subsection{电从哪里来，到哪里去}
外网是电能来源之一，负荷、光伏与电池连接在同一微网能量平衡节点。
原计划$q$是0:00承诺，调整后购电量$r$替代该承诺用于对应时段运行。
\textbf{$q$与$r$不是两股同时相加的供电量。}
正常供给不够时才另加紧急购电$e$：
\begin{equation}\label{eq:physical}
 \underbrace{r+e+G+d}_{\text{供给电量}}
 =\underbrace{L+c+u}_{\text{负荷、充电与富余处置}} .
\end{equation}
$G$为光伏电量，$L$为负荷电量，$c$为从交流母线送入电池的电量，
$d$为电池向母线送出的电量，$u$为未利用的富余电量。
储电量$s$是电池内部库存，不直接加入式\eqref{eq:physical}，
只有取出的$d$进入供电端。上述量均为kWh。

富余可来自光伏过多或已付费购电配额过多，因此不能统一限制$u\leq G$。
若另定义纯弃光$w$，才有$0\leq w\leq G$。
本模型合并记录弃光与未利用配额，不赋予售电收益。
这是一种结算抽象，不意味着存在未给定的售电市场或实际耗散设备。
若需区分付费配额和实际外网送电，可将$u$拆成未取用配额与弃光，账单仍按承诺结算。

\subsection{主线及来源}
\begin{longtable}{p{2.6cm}p{5.1cm}p{6.1cm}}
\toprule 部件 & 作用 & 来源与本题改写\\ \midrule
\endhead
随机滚动调度 & 新信息到达后重算 & Lemos-Vinasco等\cite{lemos2022}；改为6小时更新、10分钟执行。\\
提前量校准 & 区分不同远近的误差 & Schulz等\cite{schulz2021}按提前量与初始化时刻分别后处理。\\
可信度指数 & 把误差尺度变成直观评分 & 本文提出指数尺度及归一化映射，不是已有通用定律。\\
联合场景 & 保留连续短缺和变量相关性 & 借鉴概率轨迹思想\cite{lemos2022,gneiting2023}，采用成对历史日块。\\
CVaR & 权衡高额紧急费用尾部 & Rockafellar与Uryasev\cite{rockafellar2000}的辅助阈值表示。\\
决策导向选择 & 预测更准未必购电更省 & Elmachtoub与Grigas\cite{elmachtoub2022}；本题按历史调度费用选参数。\\
\bottomrule
\end{longtable}

\section{时间轴、单位与可用信息}
\subsection{144个流量值与145个库存值}
用$i$表示日期，避免与放电量$d$混淆。定义
\begin{equation}
 \Delta t=\frac16\ {\rm h},\quad T=144,\quad
 \mathcal T=\{1,\ldots,144\},\quad I_t=[(t-1)\Delta t,t\Delta t).
\end{equation}
$s_{i,t}$为第$t$时段\textbf{开始时}储电量，
$s_{i,1}$对应0:00，$s_{i,145}$对应24:00，
跨日关系为$s_{i+1,1}=s_{i,145}$。

\paragraph{时间标签歧义。}
附件1、2标签从00:10到0:00+1，而模板的时段文字整体向后10分钟。
本文按题面0:00至24:00日界，
\textbf{暂将附件标签解释为前一10分钟时段的终点}，
功率暂作该时段平均功率。00:10对应$[0{:}00,0{:}10)$，
0:00+1对应$[23{:}50,24{:}00)$。
这是需团队确认的假设，不是仅凭标签可证明的事实。
模板按行对应属于输出映射，不能据此移动真实执行日。
另一种起点解释保留核验，作图只能辅助发现错位，不能证明时间戳语义。
如有澄清，负荷、光伏、电价、预报和状态必须一起修正。

\subsection{优化时域与实际执行时域}
更新时刻$k\in\mathcal K=\{0,6,12,18\}$，单位小时。定义
\begin{equation}
 m(k)=6k+1,\quad \mathcal T_k=\{m(k),\ldots,144\},\quad
 \mathcal E_k=\{6k+1,\ldots,6(k+6)\}.
\end{equation}
每次优化当天剩余$\mathcal T_k$，只执行下一次更新前的$\mathcal E_k$。
例如6:00已执行$t=1,\ldots,36$，优化$t=37,\ldots,144$，
只执行$t=37,\ldots,72$，12:00重算。
这是\textbf{日内缩短预测时域}加终端价值，不声称每次优化完整24小时。
跨日预报可归档，目标实际值出现后再用于校准。
\begin{center}
\begin{tabular}{ccc}
\toprule 更新时刻 & 本次优化范围 & 本次执行范围\\ \midrule
0:00 & 0:00至24:00 & 0:00至6:00\\
6:00 & 6:00至24:00 & 6:00至12:00\\
12:00 & 12:00至24:00 & 12:00至18:00\\
18:00 & 18:00至24:00 & 18:00至24:00\\
\bottomrule
\end{tabular}
\end{center}

\subsection{提前量不是预报年龄}
以绝对时间记录发布时刻$a$、当前决策时刻$\tau$、目标时段终点$v$，
差值均换算为小时：
\begin{equation}\label{eq:lead}
 h=v-a,\quad A=\tau-a,\quad H=v-\tau,\quad h=A+H .
\end{equation}
$h$为预测提前量，$A$为预报年龄，$H$为距离目标还有多久。
0:00发布、针对18:00的预报，在12:00再查看，提前量仍为18小时，
不能因为距目标剩6小时就把它视为6小时提前量。
仅使用最新发布预报$a=\tau$时，才有$h=H$。
10分钟预测统一按时段终点算提前量；用中点也是备选，但校准与评价须一致。

\subsection{信息集}
$\mathcal I_{i,k}$包括当前真实库存、已结束时段的实际负荷与光伏、
此前已发布预报、已知的重复日电价、原始计划和历史参数。
不包括未来实际值及稍后才发布的预报。
附件虽有全年数据，回放仍须按时间揭示。
主校准只用当前日期以前的完整历史日；利用当天已实现误差在线修正属于后续增强，
必须另写因果更新公式。

\section{小时预报到10分钟电量}
$\widehat P^G_{i,k,j}$是$(i,k)$发布的第$j$个未来整点光伏功率，
$j=1,\ldots,24$，有效时刻为发布后$j$小时，单位kW。
$j=0$锚点用当前已获得测量；若只有10分钟记录，用刚结束时段平均值近似，
不能用即将开始时段实际值。测量不可得时用最近可见预报或历史估计。
相对时间$x\in[j,j+1]$内定义
\begin{equation}
 \widehat f_{i,k}(x)=(j+1-x)\widehat P^G_{i,k,j}
                  +(x-j)\widehat P^G_{i,k,j+1},\quad j=0,\ldots,23.
\end{equation}
对$t\in\mathcal T_k$，
\begin{align}
 x^-_{t,k}&=(t-1)\Delta t-k,\quad x^+_{t,k}=t\Delta t-k,\\
 \widehat P^G_{i,t|k}
 &=\frac1{\Delta t}\int_{x^-_{t,k}}^{x^+_{t,k}}\widehat f_{i,k}(x)\,{\rm d}x
 =\frac{\widehat f_{i,k}(x^-_{t,k})+\widehat f_{i,k}(x^+_{t,k})}{2}.
\end{align}
梯形平均成立是因10分钟网格与小时节点对齐。
若原始值表示小时平均，则改为小时内常值，作为备选解释。
插值不创造真实10分钟天气信息；历史误差须在同样重采样后与实际值比较。

第三问同样不能提前知道真实负荷。
$\mu^L_{i,t|k}$是用$\mathcal I_{i,k}$形成的负荷功率预测，
沿用问题2的历史平均、SARIMA或树模型。
光伏校准中心记作$\mu^G_{i,t|k}$。
功率场景转为电量：
\begin{equation}\label{eq:conversion}
 L^\omega_{i,t|k}=P^{L,\omega}_{i,t|k}\Delta t,\qquad
 G^\omega_{i,t|k}=P^{G,\omega}_{i,t|k}\Delta t .
\end{equation}
例如3000 kW持续10分钟为500 kWh。
费用为元/kWh乘kWh，不再乘$\Delta t$。

\section{预测提前量可信度指数}
\subsection{文献支持什么，本文新增什么}
Schulz等的《Post-processing numerical weather prediction ensembles for probabilistic
solar irradiance forecasting》\cite{schulz2021}第3.3节按提前量和初始化时刻分别
估计后处理模型，用此前若干日滚动训练。
它支持区分提前量误差，但其输入是天气集合，附件3只有单条功率预报，
不能照搬集合成员回归式。
Gneiting等的《Probabilistic solar forecasting: Benchmarks, post-processing,
verification》\cite{gneiting2023}强调校准、集中程度与适当评分，
本题借鉴这些原则而不采用其中神经网络。

“真实性”应分为：相似条件下误差范围多大，以及概率区间是否达到声称覆盖水平。
$R$描述前者，样本外概率评价检验后者。
日出、午间、夜间机制不同，不能预先断言所有情况下越近越准。

\subsection{先校正偏差}
历史样本$n$保留发布和目标时刻，原始误差
\begin{equation}
 \epsilon_n=P^{G,\rm real}_n-\widehat P^G_n\quad(\mathrm{kW}).
\end{equation}
用$g$表示少量条件组，如目标日内时刻粗分组，$\ell$表示提前量分箱。
发布时刻、目标时刻与提前量有确定关系，不视为三个独立因素随意叠加。
只有1月资料时不同时切分大量月份、小时和天气类别。
建议先共享提前量斜率，仅设少量组的偏差与尺度。
利用当前日期以前数据估计
\begin{equation}
 \widehat b_{g,\ell}
 =\frac{n_{g,\ell}\bar\epsilon_{g,\ell}+\nu_b\bar\epsilon_{\rm pool}}
        {n_{g,\ell}+\nu_b},\qquad \nu_b>0,
\end{equation}
\begin{equation}\label{eq:bias}
 \mu^G_{i,t|k}=\max\{0,\widehat P^G_{i,t|k}+\widehat b_{g,\ell}\}.
\end{equation}
$n_{g,\ell}$是不同历史日数，均值先在每天内平均，再按日等权平均；
$b,\epsilon$为kW，$n,\nu_b$无量纲。样本少则向合并组收缩，
无组样本退回合并组，无任何历史则采用冷启动。

尺度拟合使用历史滚动预测档案：
\begin{equation}\label{eq:residual}
 e_n^{\rm cal}=P^{G,\rm real}_n-\mu_n^{G,(-)}.
\end{equation}
$(-)$表示预测该历史样本时仅用更早数据。
不能把同批拟合后的训练残差当作样本外误差。

\subsection{收缩尺度估计}
$v_{g,\ell}$为校准误差方差，$v_{\rm pool}$为合并组方差，
按日等权汇总，单位$\mathrm{kW}^2$：
\begin{equation}\label{eq:shrink}
 \widetilde\sigma_{g,\ell}=
 \sqrt{\max\left\{\varepsilon_\sigma^2,
 \frac{n_{g,\ell}v_{g,\ell}+\nu_\sigma v_{\rm pool}}
      {n_{g,\ell}+\nu_\sigma}\right\}},\quad \nu_\sigma>0.
\end{equation}
$\varepsilon_\sigma>0$是数值下限，单位kW，不是物理最低误差。
残差仍有明显偏差时应先校正中心，不仅扩大方差。
确定夜间可单独设零值组；缺少位置与日历时用当时历史粗分组，
保留日出日落不确定性，不从全年测试资料倒推边界。

\subsection{本题提出的指数尺度候选}
在0至24小时内检验
\begin{equation}\label{eq:exp}
 \sigma_g(h)=\sigma_{g,0}\exp(\beta h),\quad
 \sigma_{g,0}>0,\quad\beta\geq0 .
\end{equation}
$\sigma,\sigma_{g,0}$单位kW，$h$小时，$\beta$为$\mathrm h^{-1}$，
所以$\beta h$无量纲。初期共享$\beta$，数据充足后比较分组斜率。
不向任意长期外推。
固定$s_{\rm ref}>0$为参照功率尺度，例如评价开始前历史基线白天RMSE，
用过去数据估计
\begin{equation}\label{eq:fit}
 \min_{\{a_g\},\beta\geq0}\sum_{g,\ell}n_{g,\ell}
 \left[\log\left(\frac{\widetilde\sigma_{g,\ell}}{s_{\rm ref}}\right)
       -a_g-\beta h_\ell\right]^2,\quad
 \sigma_{g,0}=s_{\rm ref}\exp(a_g).
\end{equation}
$h_\ell$为分箱代表小时数，取对数的是无量纲比值。
分箱、训练窗口与收缩强度均在过去验证期选择。
必须比较常尺度与直接分箱尺度：
若控制日内条件后仍非单调，或指数形式样本外评分差，则弃用指数形式。
\textbf{指数式是本题的可检验假设，不是文献证明的普适衰减定律。}

\subsection{归一化可信度指数及数学性质}
定义
\begin{equation}\label{eq:R}
 \boxed{R_g(h)=\frac1{1+\{\sigma_g(h)/s_{\rm ref}\}^2}
 =\frac1{1+(\sigma_{g,0}/s_{\rm ref})^2\exp(2\beta h)}},
 \quad 0<R_g(h)\leq1 .
\end{equation}
$R$无量纲，尺度越小评分越高；$R=1/2$仅表示$\sigma=s_{\rm ref}$。
同组内
\begin{equation}
 \frac{{\rm d}R_g(h)}{{\rm d}h}=-2\beta R_g(h)\{1-R_g(h)\}\leq0.
\end{equation}
跨组仍受$\sigma_{g,0}$影响，不能只按提前量排序。
也可展示相对衰减$\sigma_g(0)/\sigma_g(h)=\exp(-\beta h)$，
但它不含初始误差，不替代$R$作跨组比较。

\textbf{$R=0.8$不是80\%概率预测正确，$R=0.5$不是光伏只能发出预测的一半。}
它是人为归一化尺度评分，须同时报告$s_{\rm ref}$和独立历史日数。
参照尺度在比较期固定，高评分且样本少不代表充分证据。

\subsection{指数如何进入调度}
反解尺度
\begin{equation}\label{eq:Rscale}
 \sigma^G_{i,t|k}=s_{\rm ref}\sqrt{R_{i,t|k}^{-1}-1}.
\end{equation}
用尺度调节场景分散程度，不以$R\mu^G$替换预测中心。
否则会将无偏但波动大的预报系统性低估，还可能与CVaR重复保守。
$R$也不是场景概率，不能因不利场景可信度低就降低权重。
完整关系是
\begin{equation}
 \text{提前量与历史误差}\ \longrightarrow\ (\mu,\sigma)
 \ \longleftrightarrow\ R
 \ \longrightarrow\ \text{联合场景}
 \ \longrightarrow\ \text{购电和储能}.
\end{equation}
效果来自条件尺度校准，$R$是可解释表达，
不能将未参与决策的展示评分包装成优化贡献。

\section{由尺度产生负荷与光伏联合场景}
电池有记忆，连续一小时短缺与分散的六个10分钟误差效果不同，
因此不能独立抽144个点，也不能打乱负荷与光伏的配对。
借鉴文献\cite{lemos2022,gneiting2023}的概率轨迹思想，
使用历史日联合块，不声称完整实现其经验Copula方法。

对历史日$j$在同一发布时刻$k$的剩余路径，保存当时的预测与尺度：
\begin{equation}
 Z^X_{j,t|k}=
 \frac{P^{X,\rm real}_{j,t}-\mu^{X,(-)}_{j,t|k}}{\sigma^{X,(-)}_{j,t|k}},
 \quad X\in\{L,G\},\quad t\in\mathcal T_k .
\end{equation}
$Z$无量纲，分母正。负荷尺度可用粗分组残差，不必同样采用指数。
仅取目标已全部实现且发布时间、目标时间对齐的完整路径；
缺失、异常和近零尺度按预先规则处理。
有放回抽$M$个日块$j(\omega)$，每次同时取该日负荷与光伏：
\begin{align}
 P^{L,\omega}_{i,t|k}
 &=\max\{0,\mu^L_{i,t|k}+\sigma^L_{i,t|k}Z^L_{j(\omega),t|k}\},\\
 P^{G,\omega}_{i,t|k}
 &=\max\{0,\mu^G_{i,t|k}+\sigma^G_{i,t|k}Z^G_{j(\omega),t|k}\},\\
 \Omega&=\{1,\ldots,M\},\quad \pi_\omega=1/M,\quad\sum_\omega\pi_\omega=1.
\end{align}
然后按式\eqref{eq:conversion}转为kWh。
条件缩放与非负截断会改变最终均值、方差和相关性，须在生成场景上复核；
$\mu$称校准中心，不称截断后精确期望。
二次校准也只能用过去验证资料。
重复抽少数日期不能创造极端天气，需报告独立日数而非仅报告$M$。
不假定正态误差，不承诺覆盖所有可能天气。

\section{参数、变量与调整结算}
固定日期$i$及更新时刻$k$，以下省略日期下标。
$t|k$表示在$k$时刻为$t$制定的值，不表示除法。
\subsection{已知参数}
\begin{longtable}{p{3.2cm}p{7.0cm}p{3.6cm}}
\toprule 参数 & 定义 & 单位或范围\\ \midrule
\endhead
$p_t$ & 附件1统一时间口径下的日电价 & 元/kWh，$p_t>0$\\
$L^\omega_{t|k},G^\omega_{t|k}$ & 场景负荷与光伏电量 & kWh，非负\\
$\pi_\omega$ & 当前条件场景概率 & 非负，总和1\\
$\underline S,\overline S$ & 安全储能上下界 & 1200、10800 kWh\\
$\overline B$ & 单时段母线侧充放电上界 & $5000/6$ kWh\\
$\eta_c,\eta_d$ & 充电与放电效率 & 主口径均0.9\\
$s^{\rm obs}_{m(k)}$ & 当前真实储电量 & kWh\\
$S^{\rm tar}$ & 日末软目标，不是每日重置 & 候选6000 kWh\\
$\gamma_k$ & 日末库存缺额单位价值 & 元/kWh，非负\\
$\alpha,\lambda$ & CVaR水平与风险权重 & $0<\alpha<1,\lambda\geq0$\\
$q_t$，当$k>0$ & 0:00已承诺原始计划 & kWh，非负固定\\
\bottomrule
\end{longtable}
效率主口径均0.9，往返0.81；若90\%解释为总往返效率，两侧取$\sqrt{0.9}$比较。
不加入未给定的外网购电上限、售电价、固定调整费和退化费。

\subsection{决策变量与范围}
\begin{longtable}{p{3.1cm}p{7.0cm}p{3.7cm}}
\toprule 变量 & 含义 & 范围、单位\\ \midrule
\endhead
$q_t$，仅$k=0$ & 原始计划，求解后当天固定 & $\R_+$，kWh\\
$r_{t|k}$ & 当前版本非紧急购电量；$k=0$时等于$q_t$ & $\R_+$，kWh\\
$a^+_{t|k},a^-_{t|k}$ & 相对原始$q$的上调、下调，仅$k>0$使用 & $\R_+$，kWh\\
$c_{t|k},d_{t|k}$ & 母线侧充电与放电量 & $[0,\overline B]$，kWh\\
$s_{t|k}$ & 时段开始库存，含日末状态 & $[\underline S,\overline S]$，kWh\\
$z_{t|k}$ & 1允许充电，0允许放电 & $\{0,1\}$\\
$e^\omega_{t|k}$ & 场景紧急购电量 & $\R_+$，kWh\\
$u^\omega_{t|k}$ & 场景富余未利用电量 & $\R_+$，kWh\\
$b_k^{\rm end}$ & 日末低于软目标的库存缺额 & $\R_+$，kWh\\
$\zeta_k$ & CVaR辅助费用阈值 & $\R$，元\\
$v_\omega$ & 场景费用超阈值正偏差 & $\R_+$，元\\
\bottomrule
\end{longtable}
$\R_+=[0,\infty)$。
$q,r,a^+,a^-,c,d,s,z,b^{\rm end},\zeta$是当前场景共享决策；
$e^\omega,u^\omega,v_\omega$随场景变化。
除$s$外的调度电量均是10分钟流量。

\subsection{主结算与备选}
记$\pos{x}=\max(x,0)$，主结算
\begin{align}
 g_t(q,r)&=p_tq+1.5p_t\pos{r-q}-0.5p_t\pos{q-r}\label{eq:g}\\
 &=p_tr+0.5p_t|r-q|.
\end{align}
即最终正常购电量付费，偏离原承诺部分另付50\%；
下调1 kWh相对原计划净退回$0.5p_t$元。
计划10 kWh下调至8 kWh时结算$9p_t$元，上调至12 kWh时$13p_t$元，
仅为公式算例，不是附件结果。
若解释为原计划全付、下调另罚，则
\begin{equation}\label{eq:alt}
 g_t^{\rm alt}(q,r)=p_tq+1.5p_t\pos{r-q}+0.5p_t\pos{q-r}.
\end{equation}
主用式\eqref{eq:g}，保留式\eqref{eq:alt}核对。
多次调整按最终执行版本相对0:00计划一次结算；
若每次相对上一版增量结算，须另改费用记录。
这是题面解释，不能由文献替题面决定。

线性化为
\begin{equation}\label{eq:adjust}
 r_{t|k}=q_t+a^+_{t|k}-a^-_{t|k},\quad a^+_{t|k},a^-_{t|k}\geq0,
\end{equation}
\begin{equation}
 g_t=p_tq_t+1.5p_ta^+_{t|k}-0.5p_ta^-_{t|k}.
\end{equation}
若$a^+,a^-$同时正，可同时减少二者保持$r$不变、降低费用，
所以$p_t>0$时最优解不同时正，无需调整方向二进制变量。

\section{CVaR解释与线性化}
\subsection{风险损失对象}
定义剩余时域紧急费用
\begin{equation}\label{eq:loss}
 C_k^\omega=\sum_{t\in\mathcal T_k}5p_te^\omega_{t|k}\quad(\text{元}).
\end{equation}
紧急购电补足供需，不是未供电量。
当前正常结算不随场景改变，故对$C_k$控制尾部；
它不是未来多次重优化后整天实际账单CVaR的精确表达。

\subsection{直观与严格定义}
VaR关注分位点，CVaR还关注尾部有多严重。
$\alpha=0.95$可理解为最贵约5\%情形平均损失，
不是最坏单个场景，也不保证95\%时段不紧急购电。
离散分布在分位点有概率质量时，需取其中适当比例，
不能总写成$\E[C\mid C\geq\mathrm{VaR}_\alpha]$。
采用Rockafellar与Uryasev表示\cite{rockafellar2000}：
\begin{equation}\label{eq:cvar}
 \CVaR_\alpha(C_k)=\min_{\zeta\in\R}
 \left\{\zeta+\frac1{1-\alpha}\sum_\omega\pi_\omega\pos{C_k^\omega-\zeta}\right\}.
\end{equation}
引入
\begin{equation}\label{eq:cvarlin}
 v_\omega\geq C_k^\omega-\zeta_k,\quad v_\omega\geq0,\quad\zeta_k\in\R.
\end{equation}
目标加入$\lambda\{\zeta_k+(1-\alpha)^{-1}\sum_\omega\pi_\omega v_\omega\}$。
$\lambda=0$时删除$\zeta,v$及约束，不将未被优化的辅助值误报为风险。
损失非负时可找到非负最优阈值，这里保留一般域$\R$。

\subsection{风险偏好}
$\alpha$先比较0.90和0.95，$\lambda$选含0的少量候选，
它们不是题面参数或已验证最优值。
$C,\zeta,v$单位元，$\lambda$无量纲。
历史验证中展示实际均值与尾部权衡，降低尾部可能提高平均费。
少量历史日不能稳定识别很高分位，重复抽样不能补足信息。

\section{完整随机滚动优化模型}
\subsection{共同物理约束}
对$t\in\mathcal T_k$、$\omega\in\Omega$，
\begin{align}
 r_{t|k}+e^\omega_{t|k}+G^\omega_{t|k}+d_{t|k}
 &=L^\omega_{t|k}+c_{t|k}+u^\omega_{t|k},\label{eq:balance}\\
 s_{t+1|k}&=s_{t|k}+\eta_c c_{t|k}-d_{t|k}/\eta_d,\label{eq:soc}\\
 s_{m(k)|k}&=s^{\rm obs}_{m(k)},\label{eq:init}\\
 0\leq c_{t|k}&\leq\overline B z_{t|k},\label{eq:charge}\\
 0\leq d_{t|k}&\leq\overline B(1-z_{t|k}),\label{eq:discharge}\\
 z_{t|k}&\in\{0,1\},\quad r_{t|k},e^\omega_{t|k},u^\omega_{t|k}\geq0 .
\end{align}
包括末时刻的库存边界
\begin{equation}\label{eq:bounds}
 \underline S\leq s_{t|k}\leq\overline S,\qquad t=m(k),\ldots,145.
\end{equation}
母线取走$c$，内部只增加$\eta_cc$；输出$d$，内部减少$d/\eta_d$，
不在母线平衡重复扣效率。
互斥允许$c=d=0$，停机时$z$不唯一不影响物理。
无紧急购电上界且允许富余时，偏差可由$e,u$补救。

\subsection{日末价值与跨日库存}
仅问题1要求日首尾相同，第三问不强制每天回到6000 kWh。
为抑制缩短时域在午夜前耗尽库存，设
\begin{equation}\label{eq:terminal}
 b_k^{\rm end}\geq S^{\rm tar}-s_{145|k},\quad
 b_k^{\rm end}\geq0,\quad \Phi_k=\gamma_k b_k^{\rm end}.
\end{equation}
这是下一天库存价值近似，不是真实收费或精确动态规划价值。
$S^{\rm tar}=6000$ kWh为候选，$\gamma_k$仅用过去资料验证。
比较无终端项、正软目标和硬目标；$\gamma_k=0$时可删除$b$。
年末可另设$s_{\text{12月31日},145}=6000$ kWh作公平比较条件，
但须登记为假设而非题面强制条件。

\subsection{0:00原始计划模型}
令$r_{t|0}=q_t\geq0$，无自我调整费用：
\begin{equation}\label{eq:obj0}
 \begin{aligned}
 \min\quad&\sum_{t=1}^{144}p_tq_t+\sum_\omega\pi_\omega C_0^\omega\\
 &+\lambda\left(\zeta_0+\frac1{1-\alpha}\sum_\omega\pi_\omega v_\omega\right)
 +\gamma_0b_0^{\rm end}.
 \end{aligned}
\end{equation}
约束为$r_{t|0}=q_t\geq0$、式\eqref{eq:balance}至\eqref{eq:terminal}、
式\eqref{eq:loss}和\eqref{eq:cvarlin}。
保存$q^*$为原始承诺，只执行0:00至6:00。
当前两阶段近似未完整优化未来新预报与调整的期权价值，
不能称为精确全天多阶段最优策略。

\subsection{日内调整模型}
给定$q$及当前真实库存，只对$t\in\mathcal T_k$优化：
\begin{equation}\label{eq:objk}
 \begin{aligned}
 \min\quad&
 \sum_{t\in\mathcal T_k}(p_tq_t+1.5p_ta^+_{t|k}-0.5p_ta^-_{t|k})
 +\sum_\omega\pi_\omega C_k^\omega\\
 &+\lambda\left(\zeta_k+\frac1{1-\alpha}\sum_\omega\pi_\omega v_\omega\right)
 +\gamma_kb_k^{\rm end}.
 \end{aligned}
\end{equation}
约束为式\eqref{eq:adjust}、式\eqref{eq:balance}至\eqref{eq:terminal}、
式\eqref{eq:loss}和\eqref{eq:cvarlin}。
已执行时段不再作决策，不能追溯修改。
$\sum p_tq_t$为常数，可为求解省略，但账单解释须保留。
$q$始终是0:00版本，不改成上一版。

\subsection{模型类别与两阶段关系}
给定场景后，关系线性或已线性化；保留$z$是场景确定性等价MILP，
删除互斥、保留$0\leq c,d\leq\overline B$为LP松弛。
CVaR不使模型变成非线性规划或鲁棒优化。
LP能否替代MILP需查同时充放电和费用差。
每次第一阶段为公共$r,c,d,s,z$，第二阶段为$e^\omega,u^\omega$。
不让电池随整条未来场景任意变化，以免提前知道未来。
\textbf{整体为序贯决策，用重复两阶段近似实现随机MPC。}
334个输出日每天4次，共$334\times4=1336$次外层求解，
不含热身和校准；每次有$M$个补救块，不是独立求$1336M$个完整计划。
两次更新之间电池执行已选10分钟计划，紧急购电与富余每10分钟响应。
电池若也每10分钟自适应，必须另写因果反馈或多阶段非预见性约束。

\section{执行、跨日衔接与账单}
\subsection{只执行下一段}
对$t\in\mathcal E_k$执行$r^*,c^*,d^*$，实际值揭示后
\begin{align}
 e^{\rm real}_{i,t}
 &=\pos{L^{\rm real}_{i,t}+c^{\rm exe}_{i,t}
       -G^{\rm real}_{i,t}-d^{\rm exe}_{i,t}-r^{\rm exe}_{i,t}},\\
 u^{\rm real}_{i,t}
 &=\pos{r^{\rm exe}_{i,t}+G^{\rm real}_{i,t}+d^{\rm exe}_{i,t}
       -L^{\rm real}_{i,t}-c^{\rm exe}_{i,t}},\\
 s^{\rm real}_{i,t+1}
 &=s^{\rm real}_{i,t}+\eta_cc^{\rm exe}_{i,t}-d^{\rm exe}_{i,t}/\eta_d .
\end{align}
$e,u$不同时正且供需平衡；这是10分钟模型响应假设，不另建秒级控制。
下一更新以真实库存初始化。

\subsection{1月热身不能遗漏}
题面给1月1日6000 kWh，不是2月1日必为6000 kWh。
推荐从1月1日连续运行可执行热身，1月末状态传入2月。
初始无历史时使用官方预报与登记的简单负荷基线；
极端冷启动可行例为暂不充放电、正常购电0、由紧急购电补短缺。
这只是可行性说明而非最优建议。获得历史后建立基线与校准，
不能用1月后半月反推1月初决策。
无法产生样本外标准化残差的早期日不进入残差档案。
若不模拟1月，则2月1日6000 kWh仅为另加比较假设。
各方案相同初值后独立连续运行，不能每天重置为同一条外部库存轨迹。

\subsection{每段只结算最后有效版本}
令$\kappa(t)$为$t$前最后采用更新时间，全采用时
\begin{equation}
 \kappa(t)=6\left\lfloor\frac{t-1}{36}\right\rfloor,\quad
 r^{\rm exe}_{i,t}=r^*_{i,t|\kappa(t)}.
\end{equation}
定义$a^{+,\rm exe}_{i,t}=\pos{r^{\rm exe}_{i,t}-q_{i,t}}$、
$a^{-,\rm exe}_{i,t}=\pos{q_{i,t}-r^{\rm exe}_{i,t}}$，
真实日账单
\begin{equation}\label{eq:bill}
 F_i^{\rm real}=\sum_{t=1}^{144}
 [p_tq_{i,t}+1.5p_ta^{+,\rm exe}_{i,t}
 -0.5p_ta^{-,\rm exe}_{i,t}+5p_te^{\rm real}_{i,t}].
\end{equation}
不含CVaR惩罚与日末价值，不把四个重叠优化目标相加。
原始、各次调整、最终执行版本分别归档。
日紧急量、充放电量按10分钟电量求和，不再乘$\Delta t$。
指定区间按真实时间归属汇总；连续紧急段可合并，但不跨过中间0值。
本次仅给汇总规则，不填写附件5。

\section{可信度不等于调整价值}
\subsection{简化无储能推导}
单时段忽略电池，净需求$N\geq0$分布$F$连续、相关分位数唯一，
$p>0$，主结算，$\lambda=0$：
\begin{equation}
 \min_{r\geq0}J(r)=pq+1.5p\pos{r-q}-0.5p\pos{q-r}
                  +5p\E[\pos{N-r}].
\end{equation}
\begin{equation}
 J'(r)=\begin{cases}
 0.5p-5p\{1-F(r)\},&r<q,\\
 1.5p-5p\{1-F(r)\},&r>q.
 \end{cases}
\end{equation}
令$a=F^{-1}(0.7)$、$b=F^{-1}(0.9)$，
$F^{-1}(\tau)=\inf\{x:F(x)\geq\tau\}$，
\begin{equation}\label{eq:band}
 r^*=\begin{cases}a,&q<a,\\q,&a\leq q\leq b,\\b,&q>b.\end{cases}
\end{equation}
原计划在区间内不值得付费调整。离散分布用次梯度、分位数集合，解可不唯一。
这是按本题倍率及报童型短缺自行推导，不是论文提供的70\%与90\%常数。
电池、CVaR或不同结算会改变边界，不能逐时套入完整模型。

\subsection{误差尺度为何影响调整}
未截断位置尺度近似$N=\mu_N+\sigma_NZ$下，
\begin{equation}
 F_N^{-1}(\tau)=\mu_N+\sigma_NF_Z^{-1}(\tau).
\end{equation}
中心与尺度都改变调整区间，无需任意规定“$R$低于阈值就加购”。
净需求尺度含负荷、光伏协方差，不能仅以光伏尺度替代；
也不主张上述式在截断后原样成立。

\section{信息价值与对照实验}
\subsection{已有预报时刻}
固定0:00，对额外$\{6,12,18\}$全部8个子集$B$比较。
单独识别信息价值时，各方案在同样四个时刻重算，
仅屏蔽不在$B$的新官方预报，沿用旧预报并保留其原始提前量。
此时校准及历史日块也按旧预报的发布时刻匹配，再截取当前剩余时段，
不能误用当前发布时刻的新预报残差。
观测权限、控制机会与风险参数一致。
若连重算也取消，测量的是信息与控制更新的联合价值，应另命名。
设$J(B)=\sum_iF_i^{\rm real}(B)$，
\begin{equation}
 \Delta V(k\mid B)=J(B)-J(B\cup\{k\})\quad(\text{元}).
\end{equation}
按日配对比较，不挑有利日期。全年事后最佳子集仅是回顾性对照，
可执行选择须过去决定、未来验证。
免费信息可被忽略，在正确模型、同一分布、嵌套策略集、精确求解下，
更多信息不会提高理论最优期望费。
实际负收益可能来自校准误差、滚动近似或求解误差。
不虚构固定调整费或预报购买费。

\subsection{未提供的额外时刻}
没有3:00、9:00等新官方预报，不能宣称其实证节省率。
可在额外时刻校准旧预报再重算，但这是内部更新；
也可按提前量模型构造假设性新预报，明确标注为合成实验。
本次先完成已有四次发布的模型。
经济价值还取决于电价、库存、功率与结算，并非$R$越高越值得调整。

\subsection{一次增加一个机制}
\begin{longtable}{p{1.9cm}p{6.0cm}p{5.9cm}}
\toprule 方案 & 内容 & 验证问题\\ \midrule
\endhead
B0 & 仅0:00官方预报，日内屏蔽 & 日内官方信息价值。\\
B1 & 官方点预报、确定性MPC & 基础可执行对照。\\
B2 & 偏差校准、合并尺度、联合场景，$\lambda=0$ & 场景不确定性的作用。\\
B3 & B2改为提前量条件尺度与$R$，$\lambda=0$ & 提前量校准的作用。\\
B4 & B3增加CVaR & 尾部改善和平均费用代价。\\
尺度对照 & 常尺度、指数尺度、分箱尺度 & 指数假设是否必要。\\
选择对照 & 概率评分与历史调度费用选参 & 决策导向的增量价值。\\
\bottomrule
\end{longtable}
B0须与同控制机会方案比较；另设纯日前不修订策略则明确同时减少控制机会。
比较调度层固定预测场景，比较预测层固定调度能力。
多个机制总改善不能全部归于指数。

\section{验收标准与学习建议}
\subsection{概率预测评价}
点预测MAE、RMSE按提前量、发布时刻、昼夜分组。
覆盖目标$1-\delta$的区间为
$[\ell_n,u_n]=[F_n^{-1}(\delta/2),F_n^{-1}(1-\delta/2)]$：
\begin{align}
 {\rm PICP}&=\frac1N\sum_n\mathbf1\{\ell_n\leq y_n\leq u_n\},\\
 {\rm MPIW}&=\frac1N\sum_n(u_n-\ell_n),\\
 {\rm IS}_\delta(\ell,u;y)
 &=(u-\ell)+\frac2\delta(\ell-y)\mathbf1\{y<\ell\}
          +\frac2\delta(y-u)\mathbf1\{y>u\}.
\end{align}
$y$是实际功率，$\delta$不与CVaR的$\alpha$混用。
PICP无量纲，宽度与评分为kW，不仅扩大区间换高覆盖。
有限场景连续排序概率评分
\begin{equation}
 {\rm CRPS}_n=\sum_\omega\pi_\omega|P_n^\omega-y_n|
 -\frac12\sum_{\omega,\omega'}\pi_\omega\pi_{\omega'}
 |P_n^\omega-P_n^{\omega'}|
\end{equation}
为kW，越小越好。来源为Schulz等第3.4节\cite{schulz2021}、
Gneiting等第4节\cite{gneiting2023}。
PICP达标仍不足以证明完全校准，还应查两端分位数与条件组；
夜间零值不掩盖白天误差。这些是诊断而非有限样本概率保证。

\subsection{决策导向选参}
借鉴Smart Predict, then Optimize\cite{elmachtoub2022}，
用过去验证日$\mathcal V_i$选择少量校准参数$\theta$：
\begin{equation}
 \theta_i^*\in\arg\min_{\theta\in\Theta}
 \sum_{j\in\mathcal V_i}F_j^{\rm real}(\theta),\qquad\max\mathcal V_i<i.
\end{equation}
验证日内部预测也只用更早资料，内层选择与外层测试分离。
$\lambda$另展示均值与尾部权衡，不只按最低费便声称确定风险偏好。
这里只借鉴选择准则，不移用SPO+损失；
原文主要预测目标成本向量，本题不确定量进入约束右端，不能照搬其理论保证。

\subsection{结果评价与团队验收}
评价期2025年2月1日至12月31日共334天。
后续报告全年、月度真实费，紧急购电kWh和费用，每日紧急费经验CVaR，
调整量和次数，富余量、电池吞吐、求解时间及MILP间隙。
每日经验CVaR对实际日紧急费重新算，不加总各次优化CVaR。
按日配对并用保持日期相关性的块抽样评价变动，不把10分钟点当独立日期。
本次未执行实验，不给优劣数值。
\begin{enumerate}
 \item 时间：144段145状态，6:00第37段，午夜不漏、不重复、不整体错位。
 \item 单位：功率尺度kW，平衡电量kWh，$\beta$为$\mathrm h^{-1}$，费用元。
 \item 因果：按发布时间揭示，校准目标已实现，旧预报年龄不冒充提前量。
 \item 物理：守恒、效率、SOC、功率、互斥与跨日状态完整。
 \item 结算：$q,r,e$不混淆，下调歧义登记，每段仅一次实际结算。
 \item 指数：历史拟合与分箱对照，通过尺度进入场景，不作正确概率或均值折扣。
 \item 风险：损失对象、阈值域和参数明确，有风险中性基准，不把风险目标当账单。
 \item 证据：逐项消融、独立日数与重复场景数区分，完美信息只作同口径下界。
 \item 来源：文献与本题设计区分，指数和分位点不虚称首创。
\end{enumerate}

\subsection{最低补学要求}
依次学习提前量与滚动校准、位置尺度与收缩估计、联合残差块抽样、
CVaR阈值线性化、随机MPC非预见性、概率评分与信息价值。
检索词为\emph{forecast lead time calibration}、
\emph{location scale model shrinkage}、
\emph{joint residual block bootstrap}、
\emph{CVaR scenario linear programming}、
\emph{stochastic MPC nonanticipativity}、
\emph{probabilistic solar forecast verification}。
以现有LP、IP、随机规划知识补到能解释信息与约束关系即可，
无需先学鲁棒优化、深度学习或强化学习。

\begin{thebibliography}{99}
\bibitem{lemos2022}
Lemos-Vinasco J., Schledorn A., Pourmousavi S. A., Guericke D.
Economic evaluation of stochastic home energy management systems in a realistic rolling horizon setting.
2022, arXiv:2203.08639. \url{https://arxiv.org/abs/2203.08639}.
本地04：概率预测、随机优化、真实滚动评价。
\bibitem{schulz2021}
Schulz B., El Ayari M., Lerch S., Baran S.
Post-processing numerical weather prediction ensembles for probabilistic solar irradiance forecasting.
Solar Energy, 2021, 220: 1016--1031.
DOI: \href{https://doi.org/10.1016/j.solener.2021.03.023}{10.1016/j.solener.2021.03.023}.
\url{https://arxiv.org/abs/2101.06717}.
本地10：第3.3节提前量与发布时刻分组，第3.4节概率评分。
\bibitem{gneiting2023}
Gneiting T., Lerch S., Schulz B.
Probabilistic solar forecasting: Benchmarks, post-processing, verification.
Solar Energy, 2023, 252: 72--80.
DOI: \href{https://doi.org/10.1016/j.solener.2022.12.054}{10.1016/j.solener.2022.12.054}.
\url{https://publikationen.bibliothek.kit.edu/1000155949}.
本地11：轨迹相关性，第4节概率评价，尤其第4.4节区间评分。
\bibitem{rockafellar2000}
Rockafellar R. T., Uryasev S.
Optimization of Conditional Value-at-Risk.
Journal of Risk, 2000, 2(3): 21--41.
DOI: \href{https://doi.org/10.21314/jor.2000.038}{10.21314/jor.2000.038}.
\url{https://sites.math.washington.edu/~rtr/papers/rtr179-CVaR1.pdf}.
本地08：CVaR辅助阈值与场景表示。
\bibitem{elmachtoub2022}
Elmachtoub A. N., Grigas P.
Smart Predict, then Optimize.
Management Science, 2022, 68(1): 9--26.
DOI: \href{https://doi.org/10.1287/mnsc.2020.3922}{10.1287/mnsc.2020.3922}.
\url{https://arxiv.org/abs/1710.08005}.
本地07：借鉴决策导向，不直接移用SPO+。
\end{thebibliography}
\end{document}
